\chapter{Industrialisation et exploitation}

\section{Ce que le prototype prouve}

Le prototype démontre que la chaîne en aval de Litmus Edge peut recevoir,
contrôler, historiser et présenter les données d'une imprimante CIJ. Il ne
prouve pas encore que chaque tag simulé existe sous le même nom et dans la même
unité sur la machine du site. La mise en production ne consiste donc pas à
déplacer les conteneurs tels quels. Elle commence par une qualification du
flux réel.

Le principe d'industrialisation est de conserver les composants déjà validés
et de remplacer uniquement la source simulée et son mapping. Le courtier, le
modèle InfluxDB, les requêtes Grafana et les alertes ne doivent changer que si
les essais réels montrent une différence de sémantique, pas une simple
différence de nom.

\section{Prérequis de raccordement}

Le raccordement demande un petit ensemble d'informations, mais aucune ne peut
être supposée. Elles doivent être obtenues avant l'ouverture du flux de
production :

\begin{itemize}
  \item export des tags disponibles dans Litmus Edge, avec leur type et leur
        unité ;
  \item exemple horodaté des messages émis par le connecteur nord ;
  \item fréquence de publication et comportement lors d'une valeur inchangée ;
  \item identifiant exact de l'imprimante, de la ligne et du site ;
  \item correspondance entre les codes natifs et les états compréhensibles par
        l'utilisateur ;
  \item signal indiquant qu'une impression est attendue ;
  \item plages normales de pression, viscosité et température validées par la
        maintenance ;
  \item règles de conservation, destinataires des alertes et propriétaire de
        l'exploitation.
\end{itemize}

Sans signal de demande, l'état instantané reste visible, mais la disponibilité
ne peut pas être présentée comme un indicateur de production. Cette limite doit
être affichée explicitement plutôt que compensée par une hypothèse horaire.

\section{Table de mapping Litmus Edge}

Le mapping isole le vocabulaire du site du modèle interne de GlassInk Monitor.
Il constitue le seul endroit où les noms de tags réels sont traduits. Une table
de recette est préparée avant de modifier le flux Node-RED.

{\small
\renewcommand{\arraystretch}{1.45}
\begin{longtable}{p{0.22\textwidth}p{0.20\textwidth}p{0.13\textwidth}p{0.30\textwidth}}
\caption{Extrait de la table de mapping à compléter sur site}\\
\toprule
Champ interne & Tag Litmus & Unité attendue & Contrôle de recette \\
\midrule
\endfirsthead
\toprule
Champ interne & Tag Litmus & Unité attendue & Contrôle de recette \\
\midrule
\endhead
\nolinkurl{jet_status_code} & À renseigner & Code entier & Comparer à l'écran machine \\
\addlinespace[2pt]
\nolinkurl{printing_status_code} & À renseigner & Code entier & Lancer puis suspendre l'impression \\
\addlinespace[2pt]
\nolinkurl{line_demanding} & À renseigner & Booléen & Présenter puis retirer une pièce \\
\addlinespace[2pt]
\nolinkurl{ink_level} & À renseigner & \% & Comparer au niveau affiché \\
\addlinespace[2pt]
\nolinkurl{solvent_level} & À renseigner & \% & Comparer au niveau affiché \\
\addlinespace[2pt]
\nolinkurl{head_pressure} & À renseigner & bar & Vérifier le facteur de conversion \\
\addlinespace[2pt]
\nolinkurl{pressure_10mbar} & À renseigner & $10$ mbar & Comparer au champ précédent \\
\addlinespace[2pt]
\nolinkurl{ink_viscosity} & À renseigner & Unité machine & Observer la valeur stable \\
\addlinespace[2pt]
\nolinkurl{head_temperature} & À renseigner & $^{\circ}$C & Comparer à l'interface locale \\
\addlinespace[2pt]
\nolinkurl{print_count} & À renseigner & Compteur & Imprimer un lot connu \\
\bottomrule
\end{longtable}
}

Une conversion doit être documentée dans cette table. Par exemple, un tag de
pression exprimé en unités de 10~mbar est divisé par 100 pour obtenir des bars.
Ce calcul appartient au mapping, pas à chaque panneau Grafana. Une correction
future reste ainsi locale.

\section{Plan de mise en service}

La mise en service est découpée en six passages. Chaque passage possède une
condition de sortie ; si elle n'est pas remplie, le suivant ne commence pas.

\begin{table}[H]
\centering
\caption{Séquence de raccordement industriel}
\begin{tabularx}{\textwidth}{p{0.08\textwidth}p{0.28\textwidth}Y}
\toprule
Phase & Action & Condition de sortie \\
\midrule
P1 & Installer la VM et les volumes persistants & Contrôles de santé conformes \\
P2 & Configurer comptes, certificats, ACL et sauvegarde & Revue avec le responsable réseau \\
P3 & Recevoir le flux Litmus sur un topic de qualification & Payload capturé sans effet sur la ligne \\
P4 & Compléter le mapping et vérifier les unités & Comparaison avec l'interface de l'imprimante \\
P5 & Exécuter la recette en parallèle du suivi existant & Indicateurs acceptés par les utilisateurs \\
P6 & Activer les notifications puis transmettre l'exploitation & Procédures et responsabilités signées \\
\bottomrule
\end{tabularx}
\end{table}

Pendant P3 à P5, la plateforme reste en lecture seule. Les alertes externes
sont d'abord dirigées vers un canal de test. Elles ne sont envoyées aux
opérateurs qu'après une période d'observation suffisante pour mesurer les faux
positifs.

\section{Dimensionnement}

Le volume de la ligne ne justifie pas une architecture distribuée. Avec une
publication toutes les cinq secondes, la télémétrie produit au maximum un point
large par période. Sur une activité de 24~h/24, cinq jours par semaine, cela
représente environ :

\begin{equation}
N_{\text{télémétrie}} = 52 \times 5 \times 24 \times
\frac{3600}{5} = 4\,492\,800 \text{ points/an}
\end{equation}

Le débit maximal annoncé, inférieur à 200 pièces par heure, ajoute moins de
1,25 million d'événements de marquage par an. Les défauts et remplissages sont
beaucoup moins fréquents. Une seule VM reste donc suffisante pour le prototype
et pour une première ligne, sous réserve de mesurer l'espace disque après un
mois de données réelles.

\begin{table}[H]
\centering
\caption{Configuration initiale proposée pour une ligne}
\begin{tabularx}{\textwidth}{p{0.28\textwidth}p{0.20\textwidth}Y}
\toprule
Ressource & Valeur de départ & Raison \\
\midrule
Processeur & 4 vCPU & Marge pour InfluxDB, Grafana et les six conteneurs \\
Mémoire & 8~Go & Cache InfluxDB et rendu simultané des tableaux de bord \\
Disque & 100~Go extensibles & Historique, journaux et sauvegarde locale temporaire \\
Système & Linux 64 bits & Support stable de Docker et du pare-feu hôte \\
Réseau & Une interface supervisée & Flux limités par les règles du site \\
\bottomrule
\end{tabularx}
\end{table}

Ces valeurs sont des points de départ, pas une garantie de capacité. Après
un mois, l'exploitant relève la taille du bucket, le temps des requêtes et la
mémoire maximale. Le dimensionnement est ensuite ajusté à partir de ces
mesures.

\section{Conservation et sauvegarde}

Toutes les données n'ont pas la même valeur dans le temps. Les mesures fines
sont utiles pour comprendre les minutes qui précèdent un défaut. Plusieurs
années plus tard, une agrégation horaire suffit souvent pour une tendance. Les
marquages et défauts peuvent au contraire relever d'une obligation de
traçabilité plus longue.

\begin{table}[H]
\centering
\caption{Politique de conservation à faire valider}
\begin{tabularx}{\textwidth}{p{0.30\textwidth}p{0.22\textwidth}Y}
\toprule
Donnée & Proposition & Décision attendue \\
\midrule
Télémétrie à 5 secondes & 90 jours & Besoin d'analyse fine de la maintenance \\
Agrégats à 5 minutes & 24 mois & Comparaison des dérives et consommations \\
Événements machine & 5 ans & Historique des défauts et interventions \\
Marquages et résultats qualité & Selon la règle qualité du site & Traçabilité automobile \\
Journaux techniques & 30 jours & Diagnostic sans accumulation inutile \\
\bottomrule
\end{tabularx}
\end{table}

La sauvegarde doit inclure les volumes InfluxDB et Grafana ainsi que les
fichiers versionnés. Une copie quotidienne avec sept versions glissantes est
une base raisonnable. Une copie hebdomadaire est conservée sur un support
distinct. Le site devra fixer le RPO et le RTO ; le prototype ne peut pas les
décider à sa place. Une restauration trimestrielle sur une VM vide fournit la
preuve que les sauvegardes restent utilisables.

\section{Exploitation quotidienne}

L'exploitation ne doit pas dépendre de la personne qui a développé le
prototype. Les actions courantes sont donc limitées et associées à un rôle.

\begin{table}[H]
\centering
\caption{Répartition des responsabilités d'exploitation}
\begin{tabularx}{\textwidth}{p{0.25\textwidth}YYY}
\toprule
Action & Production & Maintenance & Administrateur \\
\midrule
Consulter l'état et acquitter une information & Oui & Oui & Oui \\
Analyser les tendances et les défauts & Lecture & Oui & Oui \\
Modifier un seuil d'alerte & Non & Proposition & Validation \\
Changer un mapping de tag & Non & Revue & Exécution \\
Gérer comptes, certificats et sauvegardes & Non & Non & Oui \\
Redémarrer un service après diagnostic & Non & Selon procédure & Oui \\
\bottomrule
\end{tabularx}
\end{table}

Un contrôle quotidien vérifie la fraîcheur du dernier point, l'état des
conteneurs et l'absence de croissance anormale des rejets. Chaque semaine, la
maintenance examine les alertes les plus fréquentes et les prévisions de
consommables. Chaque mois, l'administrateur contrôle l'espace disque, la
sauvegarde et l'expiration des certificats.

\section{Conduite d'un incident}

Lorsqu'un panneau n'affiche plus de données, le diagnostic suit le trajet de la
mesure. Il ne faut pas redémarrer tous les services sans savoir où le flux s'est
arrêté.

\begin{enumerate}
  \item confirmer l'heure du dernier point dans Grafana ;
  \item vérifier la santé d'InfluxDB et de sa source Grafana ;
  \item consulter \texttt{/health} et le journal de rejet de Node-RED ;
  \item observer le topic MQTT concerné avec un compte en lecture seule ;
  \item contrôler l'état de publication de Litmus Edge ;
  \item si le flux revient, confirmer que l'horodatage progresse et que les
        alertes se résolvent ;
  \item consigner la cause, la durée et l'action réalisée.
\end{enumerate}

Cette procédure distingue une panne de supervision d'une panne d'imprimante.
La première produit une absence de données ; la seconde continue normalement
à produire des données contenant un état de défaut.

\section{Critères d'acceptation sur site}

La recette finale doit couvrir une séquence normale et plusieurs cas limites.
Elle est menée avec un opérateur, un technicien de maintenance et le
responsable de la plateforme.

\begin{longtable}{p{0.32\textwidth}p{0.42\textwidth}p{0.14\textwidth}}
\caption{Recette fonctionnelle de mise en service}\\
\toprule
Scénario & Résultat attendu & Statut \\
\midrule
\endfirsthead
\toprule
Scénario & Résultat attendu & Statut \\
\midrule
\endhead
Démarrage puis jet prêt & Les codes et l'état opérateur suivent la machine & À tester \\
Pièce demandée et impression correcte & Classification \texttt{producing} & À tester \\
Arrêt sans commande & Aucune alerte de production & À tester \\
Arrêt pendant une demande & Alerte après temporisation & À tester \\
Variation réelle d'un niveau & Courbe, taux et prévision cohérents & À tester \\
Remplissage d'un fluide & Nouveau lot et prévision réinitialisée & À tester \\
Défaut machine connu & Code, texte et durée historisés & À tester \\
Coupure de publication & Alerte d'absence de données & À tester \\
Redémarrage de Node-RED & Reprise sans duplication visible & À tester \\
Restauration d'une sauvegarde & Tableaux et historique de nouveau accessibles & À tester \\
\bottomrule
\end{longtable}

\section{Améliorations possibles}

Une fois le flux réel stabilisé, trois prolongements deviennent pertinents.
Le premier est le passage à Sparkplug B si Litmus Edge le fournit nativement,
afin de normaliser la naissance, la mort et l'état des nœuds. Le deuxième est
une prévision enrichie par l'historique des commandes et la vitesse du jet. Le
troisième est l'ajout de plusieurs lignes avec un même modèle de dashboard et
des variables de site.

Ces extensions ne sont pas nécessaires pour prouver la solution actuelle.
Elles doivent être engagées seulement après la collecte de données réelles et
la mesure d'un besoin. Le raccordement d'une première ligne reste la priorité.
